\subsection{Equivalent measures}

PROPOSITION 10. --- Let \(\mu\) and \(\nu\) be two positive measures on T. The following conditions are equivalent:

\begin{enumerate}
\def\labelenumi{\alph{enumi})}
\item
  The locally negligible sets are the same for \(\mu\) and \(\nu\) .
\item
  The bands generated by \(\mu\) and \(\nu\) in \(\mathcal{M}(\mathrm{T})\) are identical.
\item
  One has \(\nu = g \cdot \mu\) , where \(g\) is locally \(\mu\) -integrable and \(g(t) > 0\) locally almost everywhere for \(\mu\) .
\end{enumerate}

The conditions a) and b) are equivalent by Cor. 2 of Th. 2 of No.~5. If they are satisfied, then \(\nu = g \cdot \mu\) and \(\mu = h \cdot \nu\) , where \(g\) (resp. \(h\) ) is positive and locally integrable for \(\mu\) (resp. \(\nu\) ). Therefore (No.~4, Prop. 8) \(hg\) is locally \(\mu\) -integrable and \(\mu = (hg) \cdot \mu\) . It follows (No.~3, Cor. 2 of Prop. 3) that \(hg\) is equal to 1 locally almost everywhere for \(\mu\) , so that \(g(t) > 0\) and \(h(t) = 1/g(t)\) locally almost everywhere for \(\mu\) . Conversely, suppose that \(\nu = g \cdot \mu\) with \(g(t) > 0\) locally almost everywhere for \(\mu\) ;

since \((1/g)g\) is defined locally almost everywhere and is locally \(\mu\) -integrable, 1/g is locally \(\nu\) -integrable and \((1/g)\cdot\nu=\mu\) (No.~4, Prop. 8).

DEFINITION 3. --- Two complex measures \(\theta, \theta'\) on a locally compact space T are said to be equivalent if the measures \(|\theta|\) and \(|\theta'|\) satisfy the conditions a), b), c) of Prop. 10.

For \(\theta\) and \(\theta'\) to be equivalent, it is therefore necessary and sufficient that \(|\theta|\) and \(|\theta'|\) be equivalent.

Remark. --- If \(\mu\) and \(\nu\) are two equivalent positive measures then the measurable functions defined on T, with values in any topological space G, are the same for \(\mu\) and \(\nu\) , as follows at once from Prop. 4 of No.~3.

PROPOSITION 11. --- Let \(\mu\) be a positive measure on T. If T is countable at infinity, then there exists a continuous function \(h\) such that \(h(t) > 0\) for all \(t \in \mathrm{T}\) and such that the measure \(\nu = h \cdot \mu\) (equivalent to \(\mu\) ) is bounded.

Let \((\mathrm{K}_{n})\) be a sequence of compact sets forming a covering of T and, for every n, let \(f_{n}\) be a function in \(\mathcal{K}(\mathrm{T})\) such that \(0 \leqslant f_{n} \leqslant 1\) and \(f_{n}(t) = 1\) on \(\mathrm{K}_{n}\) (Ch. III, §1, No.~2, Lemma 1). Let \((a_{n})\) be a sequence of numbers \textgreater{} 0 such that \(\sum_{n} a_{n} < +\infty\) ; the series \(h = \sum_{n} a_{n} f_{n}\) is then normally convergent in T, consequently h is a continuous function on T, such that \(h(t) > 0\) for all \(t \in \mathrm{T}\) , by construction. Setting \(\nu = h \cdot \mu\) , we then have (Prop. 3 and Ch. IV, §1, No.~3, Prop. 13)

\[
\nu^ {*} (1) = \int^ {*} h d \mu \leqslant \sum_ {n} a _ {n} \int f _ {n} d \mu .
\]

On taking for example \(a_{n} = 2^{-n}\left(\int f_{n}d\mu\right)^{-1}\) when \(\int f_n d\mu > 1\) , and \(a_{n} = 2^{-n}\) otherwise, we have \(\sum_{n}a_{n} < +\infty\) and \(\nu^{*}(1) < +\infty\) , which proves the proposition.

PROPOSITION 12. --- Let \((\mu_n)\) be a sequence of bounded positive measures on T; there exists a bounded positive measure \(\mu\) on T such that the relation \(\mu^*(\mathrm{N}) = 0\) is equivalent to \(\langle \mu_n^*(\mathrm{N}) = 0 \text{ for every } n \rangle\) ; each of the measures \(\mu_n\) has base \(\mu\) . Moreover, if \(\mu'\) is a second positive measure on T having this property, then \(\mu\) and \(\mu'\) are equivalent.

The last part of the statement follows at once from Def. 3. To prove the existence of \(\mu\) , we can restrict ourselves to the case that \(\mu_n \neq 0\) for every \(n\) ; the family of measures \(\mu_n / 2^n \| \mu_n \|\) is then summable in \(\mathcal{M}(\mathrm{T})\) , and its sum \(\mu\) is such that \(\| \mu \| \leqslant 1\) . Moreover, since \(\mu_n \leqslant 2^n \| \mu_n \| \cdot \mu\) , the relation \(\mu(\mathrm{N}) = 0\) implies that \(\mu_n(\mathrm{N}) = 0\) for all \(n\) ; conversely, if \(\mathrm{N}\) is a set that is negligible for all the \(\mu_n\) , then it is locally negligible for \(\mu\) (§2, No.~2, Cor. 2 of Prop. 1), hence is \(\mu\) -negligible since \(\mu\) is bounded (§1, No.~2, Cor. 2 of Prop. 7).

